\chapter*{\SA{निर्णयमानचित्रम्} — decision map}
\addcontentsline{toc}{chapter}{निर्णयमानचित्रम्}
\begin{bardown}{B0 — क्तनिर्णयस्य मुख्यप्रवाहः}
\BRow{\textbf{Input gate}}{Identify the exact Dhātupāṭha entry, not merely a dictionary spelling. Record gaṇa, indicatory markers, meaning, pada information, preverb, and intended syntactic interpretation.}
\BRow{\textbf{Affix gate}}{Invoke \RuleRef{AS-3-2-102}; register क्त as निष्ठा through \RuleRef{AS-1-1-26}.}
\BRow{\textbf{Voice gate}}{Test ordinary कर्मणि/resultative interpretation; then test \RuleRef{AS-3-4-72}.}
\BRow{\textbf{iṭ gate}}{Apply the ārdhadhātuka iṭ system beginning with \RuleRef{AS-7-2-35}; do not decide seṭ/aniṭ/veṭ from gaṇa alone.}
\BRow{\textbf{Stem gate}}{Apply root-specific and suffix-conditioned operations in grammatical order.}
\BRow{\textbf{Niṣṭhā gate}}{Search 8.2.42ff for substitution, optionality, prohibition, and meaning-conditioned forms.}
\BRow{\textbf{Sandhi gate}}{Apply only sandhi and tripādī operations actually triggered.}
\BRow{\textbf{Audit gate}}{Return surface form, derivation path, competing forms, interpretation, source status, and unresolved commentarial question.}
\end{bardown}
